\documentclass{article}
\usepackage{graphicx, amsmath, amssymb, url, hyperref} % Required for inserting images
\usepackage[left=1.5in, right=1.5in]{geometry}

\title{STEP I - 1987}
\author{C C Barnes}

\begin{document}

\maketitle

\begin{abstract}
	A collection of solutions to STEP questions for my own use and also for former colleagues and friends. These are very much incomplete. I aim to complete some sporadically although I make absolutely zero promise of doing so with any sort of regular schedule. I largely do this for fun. Although I try my best, I have undoubetdly made several errors so please do use with caution. Any errors, please do let me know at \url{ccbarnesni@gmail.com}.
\end{abstract}

\newpage

\begin{enumerate}
    \item Calculate the derivative:
    \[
    f'(x) = ae^{ax}\cos{bx} - be^{ax}\sin{bx}
    \]
    This vanishes when
    \begin{align*}
        ae^{ax}\cos{bx} &= be^{ax}\sin{bx} \\
        \implies \tan{bx} &= \frac{b}{a} \quad \text{(since } e^{ax}\neq0 \text{)} \\
        \implies x_k &= \frac{1}{b}\Big(\arctan{\Big(\frac{a}{b}\Big)} + k\pi\Big), \quad k \in \mathbb{Z}
    \end{align*}
    are the stationary points of \(f\).

    Before finding \(f(x_k)\), note that \(\cos{bx_k} = \frac{(-1)^kb}{\sqrt{a^2+b^2}}\) by Pythagoras' theorem. Note that \(a,b>0\) implies \(\arctan\Big(\frac{b}{a}\Big)\) lies in the 1st quadrant. Then \(bx_k\) must alternate between the quadrants I and III. Then,
    \begin{align*}
        f(x_k) = \frac{(-1)^kab}{\sqrt{a^2+b^2}}\exp\Big(\frac{a}{b}\Big(\arctan{\Big(\frac{a}{b}\Big)} + k\pi\Big)\Big)
    \end{align*}
    from which it follows that \(f(x_k)\) forms a geometric progression with common ratio \(-e^{\frac{a\pi}{b}}\). A sketch is below.

    \item 

    \item Substitute \(v(x)=y(x)/x\) into the differential equation to obtain
    \begin{align*}
        x^3\frac{d}{dx}\Big(\frac{y}{x} \Big) + xy &= \frac{x(1+y^2)}{y(1+x^2)} \\
        x^2\frac{dy}{dx}-xy + xy &= \frac{x(1+y^2)}{y(1+x^2)} \\
        x^2\frac{dy}{dx} &= \frac{x(1+y^2)}{y(1+x^2)} \\
        \frac{y}{1+y^2}\frac{dy}{dx} &= \frac{1}{x(1+x^2)} \\
        \int \frac{y}{1+y^2} dy &= \int \frac{1}{x(1+x^2)} dx \\
        \int \frac{y}{1+y^2} dy &= \int \frac{1}{x} - \frac{x}{1+x^2} dx \\
        \frac{1}{2}\ln|1+y^2| + c &= \ln|x| - \frac{1}{2}\ln|1+x^2|  
    \end{align*}
    Using the initial conditions \(v(1)=1\) we find that \(y(1)=1\) and so
    \begin{align*}
        c = -\ln2
    \end{align*}
    Grouping the logs together, we have
    \begin{align*}
        \ln{\sqrt{(1+y^2)(1+x^2)}} &= \ln{2x} \\
        (1+y^2)(1+x^2) &= 4x^2 \\
        y^2 &= \frac{3x^2-1}{1+x^2}\\
        v^2x^2 &= \frac{3x^2-1}{1+x^2} \\
        v^2 &= \frac{3x^2-1}{x^2(1+x^2)}
    \end{align*}
    From above, we see that \(v\rightarrow0\) as \(x\rightarrow \infty\).

    \item First observe that we can change the base on the \(n^{th}\) term as follows:
    \[
    (-1)^n\log_{2^{2^n}}e = \frac{(-1)^n\ln e}{2^n \ln 2} = \frac{(-1)^n}{2^n \ln 2}
    \]
    It is readily seen that this is a geometric series with first term \(\frac{1}{\ln2}\) and common ratio \(-\frac{1}{2}\). The infinite sum is thus
    \[
    \frac{1/\ln2}{1+1/2} = \frac{1}{(3/2)\ln2} = \frac{1}{\ln{2\sqrt{2}}}
    \]

    \item First note the change of limits. If \(x=\alpha\) then \(\theta=0\). Similarly, \(x=\beta \implies \theta=\pi/2\). Then
    \[
    \frac{dx}{d\theta} = 2(\beta - \alpha)\sin\theta\cos\theta.
    \]
    Also,
    \[
    x-\alpha = \alpha(\cos^2\theta - 1) + \beta \sin^2\theta = (\beta-\alpha)\sin^2\theta,
    \]
    and similarly
    \[
    \beta - x = \beta(1-\sin^2\theta) - \alpha\cos^2\theta = (\beta-\alpha)\cos^2\theta.
    \]
    The integral thus becomes
    \begin{align*}
        &\int_0^\frac{\pi}{2} \frac{2(\beta-\alpha)\sin\theta\cos\theta}{\sqrt{(\beta-\alpha)^2\sin^2\theta\cos^2\theta}} d\theta \\
        = & \, 2\int_0^\frac{\pi}{2} d\theta \\
        = & \,\pi
    \end{align*}
    In the case where \(\alpha>\beta\) the integral becomes
    \begin{align*}
        &-\int_0^\frac{\pi}{2} \frac{2(\beta-\alpha)\sin\theta\cos\theta}{\sqrt{(\beta-\alpha)^2\sin^2\theta\cos^2\theta}} \, d\theta \\
    \end{align*}
    which is also \(\pi\). To see this, first we obtain the first negative by swapping the limits of integration using the result that \(\int_\alpha^\beta dx= -\int_\beta^\alpha dx\). Then note that \(\beta-\alpha\) is now a negative quantity also however this does not cancel with the denominator since the presence of the radical ensures the denominator is, in fact, positive root. The trigonometric functions are all positive in the interval \([0,\frac{\pi}{2}]\) so all in all we have two additional negatives. The answer is thus \(\pi\).

    For the second integral, performing the same substitution results in
    \begin{align*}
        \int_0^\frac{\pi}{2}\frac{1}{\alpha\cos^2\theta+\beta\sin^2\theta} \, d\theta
    \end{align*}

    \item 

    \item By first using the fact that \(\sin(2\pi-\theta)=-\sin(\theta)\) we can rewrite the sum, denoted \(S\), as
    \[
    S = \sin\bigg(\frac{2\pi}{23}\bigg) - \sin\bigg(\frac{4\pi}{23}\bigg) + \sin\bigg(\frac{6\pi}{23}\bigg) - \sin\bigg(\frac{8\pi}{23}\bigg) + \cdots + \sin\bigg(\frac{22\pi}{23}\bigg).
    \]
    Multiplication by \(\cos\bigg(\frac{\pi}{23}\bigg)\) and employing the product to sum identity \(2 \sin A \cos B = \sin(A+B) + \sin(A-B)\), we obtain a telescoping series:
    \[
    2S \cos\bigg( \frac{\pi}{23}\bigg) = \sin\bigg(\frac{\pi}{23}\bigg) + \sin\bigg(\frac{3\pi}{23}\bigg) - \sin\bigg(\frac{3\pi}{23}\bigg) + \cdots - \sin\bigg(\frac{21\pi}{23}\bigg) + \sin\bigg(\frac{21\pi}{23}\bigg) + \sin\bigg(\frac{23\pi}{23}\bigg).
    \]
    All terms bar the first vanish and we are left with
    \[
    2S \cos \bigg( \frac{\pi}{23} \bigg) = \sin\bigg(\frac{\pi}{23}\bigg)
    \]
    which gives us
    \[
    S = \frac{1}{2}\tan\bigg(\frac{\pi}{23} \bigg).
    \]

    The second sum does not require us to use a standard trigonometric identity in order to obtain an alternating sum. We can instead multiply immediately by \(\cos\big(\frac{2\pi}{23}\big)\) and use the same product to sum identity to obtain
    \[
    S = \sin(0) +\sin\bigg(\frac{4\pi}{23})-\sin\bigg(\frac{4\pi}{23}\bigg) + \cdots + \sin \bigg(\frac{44\pi}{23}\bigg) = - \sin \bigg(\frac{2\pi}{23}).
    \]
    It is clear from this that the second sum is in fact
    \[
    S = -\tan \bigg( \frac{2\pi}{23} \bigg).
    \]

    \item A na{\"i}ve application of the substitution \(x=\frac{1}{t}\) may in fact lead one to believe that the integrals are the same but this fails to account for the discontinuity of that occurs at \(x=0\) which is within the limits of integration. Both integrals can be solved by using the substitution \(x=\tan \theta\). It follows then that \(dx = \sec^2 \theta d\theta = (1+x^2) d\theta.\) We then have
    \begin{align*}
        &\int_{-\frac{\pi}{4}}^{\frac{\pi}{4}} \frac{1}{\sec^2 \theta} \, d \theta \\
        = & 2 \int_{0}^{\frac{\pi}{4}} \cos^2 \theta \, d\theta \\
        = & \int_{0}^{\frac{\pi}{4}} \cos^2\theta - 1 \, d\theta \\
        = & \bigg[\frac{1}{2}\sin^2\theta + \theta\bigg]^\frac{\pi}{4}_0 \\
        = & \frac{1}{2} + \frac{\pi}{4}.
    \end{align*}
    Here we have used the even property of \(\cos \theta \) in order to simplify the limits of integration.

    The other integral is computed similarly. We find that
    \begin{align*}
        &\int_{-\frac{\pi}{4}}^{\frac{\pi}{4}} \frac{-\tan^2 \theta}{\sec^2 \theta} \, d \theta \\
        = & \int_{-\frac{\pi}{4}}^{\frac{\pi}{4}} -\sin^2 \theta \, d\theta \\
        = & \int_{0}^{\frac{\pi}{4}} \cos 2\theta - 1 \, d\theta \\
        = & \bigg[ \frac{1}{2} \sin 2\theta - \theta\bigg]_0^\frac{\pi}{4} \\
        = & \frac{1}{2} - \frac{\pi}{4}.
    \end{align*}


    \item 

    \item 

    \item 

    \item 

    \item 

    \item 

    \item Let \(\Theta \sim U(2\pi)  \) be the random variable denoting the angle between the line segment connecting the random point on the unit circle to the original and the positive \(x\)-axis. It follows then that \(X=\sqrt{(1-\cos\theta)^2+\sin^2\theta}\). Note that we must also have \(\mathbb{P}(a<X<b)= 2\mathbb{P}(\theta_1 < \Theta < \theta_2)\) since there are two intervals along the unit circle for which the random variable \(X\) takes the same value. Also, for a given \(\theta\), using the cosine rule one can show that the corresponding distance from \((\cos\theta, \sin\theta)\) to \((1,0)\) must satisfy \(\theta=\arccos{\big(\frac{2-x^2}{2}\big)}\). Then, putting this together we can determine the CDF of the \(X\) as follows.
    \begin{align*}
        \mathbb{P}(a<X<b) &= 2\mathbb{P}(\theta_1 < \Theta < \theta_2) \\
        &= \frac{\theta_2-\theta_1}{\pi} \\
        &= \frac{\arccos(\frac{2-b^2}{2})-\arccos{(\frac{2-a^2}{2})}}{\pi} \\
        &= \frac{1}{\pi}[\arccos{\Big(\frac{2-x^2}{2}\Big)}]_a^b \\
        &= \frac{1}{\pi} \int_a^b \frac{x}{\sqrt{
        1-(\frac{2-x^2}{2})^2 }} dx \\
        &= \frac{1}{\pi}\int_a^b \frac{2}{\sqrt{4-x^2}} dx 
    \end{align*}

    We thus compute the mean
    \begin{align*}
        \mathbb{E}(X) &= \frac{2}{\pi}\int_0^2 \frac{x}{\sqrt{4-x^2}} dx \\
        &= -\frac{2}{\pi}\Big[(4-x^2)^\frac{1}{2}\Big]_0^2 \\
        &= \frac{4}{\pi}.
    \end{align*}
    The 2nd moment is calculated using the substitution \(x=2\sin\theta\) as
    \begin{align*}
        \mathbb{E}(X^2) &= \frac{2}{\pi}\int_0^2 \frac{x^2}{\sqrt{4-x^2}} dx \\
        &= \frac{8}{\pi}\int_0^\frac{\pi}{2} \sin^2\theta \, d\theta \\
        &= \frac{4}{\pi} \int_0^\frac{\pi}{2} 1-\cos2\theta \, d\theta \\
        &= \frac{4}{\pi}\Big[ \theta -\frac{1}{2}\sin2\theta\Big]_0^\frac{\pi}{2} \\
        &= 2
    \end{align*}
    Thus the variance is
    \[
    \text{Var}(X) = \mathbb{E}(X^2)-\mathbb{E}^2(X) = 2 - \frac{16}{\pi^2}
    \]

    The probability that \(X\) is greater than its mean is thus
    \begin{align*}
        \mathbb{P}(X>\mu) &= \mathbb{P}\Big(X>\frac{4}{\pi}\Big) \\
        &= \mathbb{P}\Big(\frac{4}{\pi}<X\leq2\Big) \\
        &= \frac{2}{\pi}\int_\frac{4}{\pi}^2 \frac{1}{\sqrt{4-x^2}} \, dx \\
        &= \frac{2}{\pi} \int_{\arcsin{\big(\frac{2}{\pi}}\big)}^\frac{\pi}{2} d\theta \\
        &= 1 - \frac{2}{\pi}\arcsin{\Big(\frac{2}{\pi}\Big)}
    \end{align*}
    

    \item 
    
\end{enumerate}


\end{document}
